% arara: pdflatex: { files: [latexindent]}
\subsubsection{The remaining code blocks}
 Referencing the different types of code blocks in \vref{tab:code-blocks}, we have a few
 code blocks yet to cover; these are very similar to the \texttt{commands} code block
 type covered comprehensively in \vref{subsubsec:commands-arguments}, but a small
 discussion defining these remaining code blocks is necessary.

 \paragraph{keyEqualsValuesBracesBrackets}
  \texttt{latexindent.pl} defines this type of code block by the following criteria:
  \begin{itemize}
   \item it has a name made up of the characters detailed in \vref{tab:code-blocks};
   \item then an $=$ symbol;
   \item then at least one set of curly braces or square brackets (comments and line
         breaks allowed throughout).
  \end{itemize}
  See the \texttt{keyEqualsValuesBracesBrackets: name} field of the fine tuning section
  in \vref{lst:fineTuning} \announce{2019-07-13}{fine tuning:
  keyEqualsValuesBracesBrackets}%

  \begin{example}
  An example is shown in \cref{lst:pgfkeysbefore}, with the default output given in
  \cref{lst:pgfkeys1:default}.

  \begin{cmhtcbraster}
   \cmhlistingsfromfile{demonstrations/pgfkeys1.tex}{\texttt{pgfkeys1.tex}}{lst:pgfkeysbefore}
   \cmhlistingsfromfile[showtabs=true]{demonstrations/pgfkeys1-default.tex}{\texttt{pgfkeys1.tex} default output}{lst:pgfkeys1:default}
  \end{cmhtcbraster}

  In \cref{lst:pgfkeys1:default}, note that the maximum indentation is three tabs, and
  these come from:
  \begin{itemize}
   \item the \lstinline!\pgfkeys! command's mandatory argument;
   \item the \lstinline!start coordinate/.initial! key's mandatory argument;
   \item the \lstinline!start coordinate/.initial! key's body, which is defined as any
         lines following the name of the key that include its arguments. This is the part
         controlled by the \emph{body} field for \texttt{noAdditionalIndent} and friends
         from \cpageref{sec:noadd-indent-rules}.
  \end{itemize}
  \end{example}

 \paragraph{namedGroupingBracesBrackets} This type of code block is mostly motivated by
  tikz-based code; we define this code block as follows:
  \begin{itemize}
   \item the name may contain the characters detailed in \vref{tab:code-blocks};
   \item then at least one set of curly braces or square brackets (comments and line
         breaks allowed throughout).
  \end{itemize}
  See the \texttt{NamedGroupingBracesBrackets: name} field of the fine tuning section in
  \vref{lst:fineTuning} \announce{2019-07-13}{fine tuning: namedGroupingBracesBrackets}%

  \begin{example}
  A simple example is given in \cref{lst:child1}, with default output in
  \cref{lst:child1:default}.

  \begin{cmhtcbraster}
   \cmhlistingsfromfile{demonstrations/child1.tex}{\texttt{child1.tex}}{lst:child1}
   \cmhlistingsfromfile[showtabs=true]{demonstrations/child1-default.tex}{\texttt{child1.tex} default output}{lst:child1:default}
  \end{cmhtcbraster}

  In particular, \texttt{latexindent.pl} considers \texttt{child} and \texttt{parent} to
  be \texttt{namedGroupingBracesBrackets}\footnote{ You may like to verify this by using
  the \texttt{-t} option and checking \texttt{indent.log}! }. Referencing
  \cref{lst:child1:default}, note that the maximum indentation is three tabs, and these
  come from:
  \begin{itemize}
   \item the \lstinline!child!'s mandatory argument;
   \item the \lstinline!child!'s body, which is defined as any lines following the name
         of the \texttt{namedGroupingBracesBrackets} that include its arguments. This is
         the part controlled by the \emph{body} field for \texttt{noAdditionalIndent} and
         friends from \cpageref{sec:noadd-indent-rules};
   \item the \lstinline!parent!'s body, which is defined as any lines following the name
         that include its arguments.
  \end{itemize}
  \end{example}

  \begin{example}
  Consider the file given in \cref{lst:named1}, together with its default output using
  the command

  \begin{commandshell}
latexindent.pl named1.tex 
\end{commandshell}

  is given in \cref{lst:named1-default}.

  \begin{cmhtcbraster}[raster column skip=.01\linewidth]
   \cmhlistingsfromfile{demonstrations/named1.tex}{\texttt{named1.tex}}{lst:named1}
   \cmhlistingsfromfile{demonstrations/named1-default.tex}{\texttt{named1.tex} default}{lst:named1-default}
  \end{cmhtcbraster}

  \end{example}

  \begin{example}
  Consider the file given in \cref{lst:finetuning2}, together with its default output
  using the command

  \begin{commandshell}
latexindent.pl finetuning2.tex 
\end{commandshell}

  is given in \cref{lst:finetuning2-default}.

  \begin{cmhtcbraster}[raster column skip=.01\linewidth,
    raster left skip=-3.75cm,
    raster right skip=0cm,]
   \cmhlistingsfromfile{demonstrations/finetuning2.tex}{\texttt{finetuning2.tex}}{lst:finetuning2}
   \cmhlistingsfromfile{demonstrations/finetuning2-default.tex}{\texttt{finetuning2.tex} default}{lst:finetuning2-default}
  \end{cmhtcbraster}

  \end{example}

  \begin{example}

  Starting with the file in \cref{lst:bib1} and running the command \index{bibliography
  files} \index{regular expressions!delimiterRegEx} \index{regular expressions!capturing
  parenthesis} \index{regular expressions!ampersand alignment}
  \index{delimiters!delimiterRegEx}

  \begin{commandshell}
latexindent.pl bib1.tex -o=+-mod1 
   \end{commandshell}

  gives the output in \cref{lst:bib1-mod1}.

  \begin{widepage}
   \begin{cmhtcbraster}[raster column skip=.01\linewidth]
    \cmhlistingsfromfile{demonstrations/bib1.bib}{\texttt{bib1.bib}}{lst:bib1}
    \cmhlistingsfromfile{demonstrations/bib1-mod1.bib}{\texttt{bib1-mod1.bib}}{lst:bib1-mod1}
   \end{cmhtcbraster}
  \end{widepage}

  Let's assume that we would like to format the output so as to align the \texttt{=}
  symbols. Using the settings in \cref{lst:bibsettings1} and running the command

  \begin{commandshell}
latexindent.pl bib1.bib -l bibsettings1.yaml -o=+-mod2 
     \end{commandshell}

  gives the output in \cref{lst:bib1-mod2}.

  \begin{widepage}
   \begin{cmhtcbraster}[raster column skip=.1\linewidth]
    \cmhlistingsfromfile{demonstrations/bib1-mod2.bib}{\texttt{bib1.bib} using \cref{lst:bibsettings1}}{lst:bib1-mod2}
    \cmhlistingsfromfile[style=yaml-LST]{demonstrations/bibsettings1.yaml}[yaml-TCB]{\texttt{bibsettings1.yaml}}{lst:bibsettings1}
   \end{cmhtcbraster}
  \end{widepage}
  we have populated the \texttt{lookForAlignDelims} field with the \texttt{online}
  command, and have used the \texttt{delimiterRegEx}, discussed in
  \vref{sec:delimiter-reg-ex}.
  \end{example}

  \begin{example}
  We can build upon \cref{lst:bibsettings1} for slightly more complicated bibliography
  files.

  Starting with the file in \cref{lst:bib2} and running the command

  \begin{commandshell}
latexindent.pl bib2.bib -l bibsettings1.yaml -o=+-mod1 
   \end{commandshell}

  gives the output in \cref{lst:bib2-mod1}.

  \begin{widepage}
   \cmhlistingsfromfile{demonstrations/bib2.bib}{\texttt{bib2.bib}}{lst:bib2}
   \cmhlistingsfromfile{demonstrations/bib2-mod1.bib}{\texttt{bib2-mod1.bib}}{lst:bib2-mod1}
  \end{widepage}

  The output in \cref{lst:bib2-mod1} is not ideal, as the \texttt{=} symbol within the
  url field has been incorrectly used as an alignment delimiter.

  We address this by tweaking the \texttt{delimiterRegEx} field in
  \cref{lst:bibsettings2}.

  \cmhlistingsfromfile[style=yaml-LST]{demonstrations/bibsettings2.yaml}[yaml-TCB]{\texttt{bibsettings2.yaml}}{lst:bibsettings2}

  Upon running the command

  \begin{commandshell}
latexindent.pl bib2.bib -l bibsettings1.yaml,bibsettings2.yaml -o=+-mod2 
         \end{commandshell}

  we receive the \emph{desired} output in \cref{lst:bib2-mod2}.

  \cmhlistingsfromfile{demonstrations/bib2-mod2.bib}{\texttt{bib2-mod2.bib}}{lst:bib2-mod2}

  With reference to \cref{lst:bibsettings2} we note that the \texttt{delimiterRegEx} has
  been adjusted so that \texttt{=} symbols are used as the delimiter, but only when they
  are \emph{not preceded} by either \texttt{v} or \texttt{spfreload}.
  \end{example}

 \paragraph{UnNamedGroupingBracesBrackets} occur in a variety of situations;
  specifically, we define this type of code block as satisfying the following
  criteria:
  \begin{itemize}
   \item it isn't a \texttt{command}, \texttt{keyEqualsValuesBracesBrackets} or
         \texttt{namedGroupingBracesBrackets}
   \item and it has at least one set of curly braces or square brackets (comments and
         line breaks allowed throughout).
  \end{itemize}

  \begin{example}
  An example is shown in \cref{lst:psforeach1} with default output give in
  \cref{lst:psforeach:default}.

  \begin{cmhtcbraster}
   \cmhlistingsfromfile{demonstrations/psforeach1.tex}{\texttt{psforeach1.tex}}{lst:psforeach1}
   \cmhlistingsfromfile[showtabs=true]{demonstrations/psforeach1-default.tex}{\texttt{psforeach1.tex} default output}{lst:psforeach:default}
  \end{cmhtcbraster}

  Referencing \cref{lst:psforeach:default}, there are \emph{two} sets of unnamed braces.
  Note also that the maximum value of indentation is two tabs, and these come from:
  \begin{itemize}
   \item the \lstinline!\psforeach! command's mandatory argument;
   \item the \emph{first} un-named braces mandatory argument.
  \end{itemize}
  \end{example}

  Users wishing to customise the mandatory and/or optional arguments on a \emph{per-name}
  basis for the \texttt{UnNamedGroupingBracesBrackets} should use
  \texttt{always-un-named}.

  \begin{example}
  Starting with the file in \cref{lst:unnamed1} we receive the default output in
  \cref{lst:unnamed1:mod0}.

  \begin{cmhtcbraster}
   \cmhlistingsfromfile*{demonstrations/unnamed1.tex}{\texttt{unnamed1.tex}}{lst:unnamed1}
   \cmhlistingsfromfile*[showtabs=true]{demonstrations/unnamed1-mod0.tex}{\texttt{unnamed1.tex} default output}{lst:unnamed1:mod0}
  \end{cmhtcbraster}

  We can customise the output using, for example, \cref{lst:unnamed1:yaml} and running

  \begin{commandshell}
latexindent.pl unnamed1 -l unnamed1.yaml unnamed1
         \end{commandshell}

  gives the output in \cref{lst:unnamed1:mod0}

  \begin{cmhtcbraster}
   \cmhlistingsfromfile*[style=yaml-LST]{demonstrations/unnamed1.yaml}[yaml-TCB]{\texttt{unnamed1.yaml}}{lst:unnamed1:yaml}
   \cmhlistingsfromfile*[showspaces=true]{demonstrations/unnamed1-mod1.tex}{\texttt{unnamed1.tex} mod1 output}{lst:unnamed1:mod1}
  \end{cmhtcbraster}
  \end{example}

\subsubsection{Summary}
 \index{indentation!summary}
 Having considered all of the different types of code blocks, the functions of the fields
 given in \cref{lst:noAdditionalIndentGlobal,lst:indentRulesGlobal} should now make sense.
 \index{specialBeginEnd!noAdditionalIndentGlobal}
 \index{specialBeginEnd!indentRulesGlobal}

 \begin{widepage}
  \begin{minipage}{.45\linewidth}
   \cmhlistingsfromfile[style=noAdditionalIndentGlobal]{../defaultSettings.yaml}[before=\centering,yaml-TCB]{\texttt{noAdditionalIndentGlobal}}{lst:noAdditionalIndentGlobal}
  \end{minipage}%
  \hfill
  \begin{minipage}{.49\linewidth}
   \cmhlistingsfromfile[style=indentRulesGlobal]{../defaultSettings.yaml}[before=\centering,yaml-TCB]{\texttt{indentRulesGlobal}}{lst:indentRulesGlobal}
  \end{minipage}%
 \end{widepage}
